\documentclass[11pt,a4paper]{article}

\usepackage[T1]{fontenc}
\usepackage[utf8]{inputenc}
\usepackage{mathpazo}
\usepackage{amsmath,amssymb,amsthm}
\usepackage[margin=2.6cm]{geometry}
\usepackage{booktabs}
\usepackage{enumitem}
\usepackage{microtype}
\usepackage{xcolor}
\usepackage[font=small,labelfont=bf]{caption}
\usepackage{listings}
\usepackage{url}
\usepackage{fancyhdr}
\usepackage[hidelinks,breaklinks,
  pdftitle={Coincidence Classes in the ECDSA State Space: A Zero-Yield Theorem},
  pdfauthor={Andriy Polishchuk},
  pdfsubject={Negative result on coincidence classes in ECDSA; secp256k1},
  pdfkeywords={ECDSA, secp256k1, lattice-point counting, negative result, hidden number problem}
]{hyperref}

\definecolor{errcol}{RGB}{138,44,44}
\definecolor{acccol}{RGB}{20,82,75}

\theoremstyle{plain}
\newtheorem{theorem}{Theorem}
\newtheorem{lemma}[theorem]{Lemma}
\newtheorem{proposition}[theorem]{Proposition}
\newtheorem{corollary}[theorem]{Corollary}
\theoremstyle{definition}
\newtheorem{definition}[theorem]{Definition}
\newtheorem{assumption}{Assumption}
\renewcommand{\theassumption}{\Alph{assumption}}
\theoremstyle{remark}
\newtheorem{remark}[theorem]{Remark}
\newtheorem{erratum}{Erratum}

\newcommand{\Zs}{\mathbb{Z}_n^{*}}
\newcommand{\Zz}{\mathbb{Z}}
\newcommand{\Qu}{\mathcal{Q}}
\newcommand{\Ex}{\mathbb{E}}
\newcommand{\Pro}{\Pr}
\newcommand{\szk}{s_{zk}}
\newcommand{\srxk}{s_{rxk}}
\newcommand{\szr}{s_{zr}}
\newcommand{\clA}{\mathrm{A}}
\newcommand{\clB}{\mathrm{B}}
\newcommand{\clE}{\mathrm{E}}
\newcommand{\clC}{\mathrm{C}}
\newcommand{\clD}{\mathrm{D}}

\lstset{
  basicstyle=\ttfamily\footnotesize,
  breaklines=true, frame=single, framesep=5pt,
  rulecolor=\color{black!30}, columns=fullflexible,
  keywordstyle=\color{acccol}, commentstyle=\color{black!55}\itshape,
  language=Python, showstringspaces=false, numbers=none
}

\title{\bfseries Coincidence Classes in the ECDSA State Space:\\[2pt]
A Zero-Yield Theorem, with Exact Frequencies\\ for Classes A, B and E}

\author{%
  \textsc{Andriy Polishchuk}\thanks{Independent researcher. CRYPTON Systems Lab is an
  independent research group and not an institutional affiliation.}\\[2pt]
  \small CRYPTON Systems Lab\\
  \small\texttt{andriy.polishchuk.a@gmail.com}
}
\date{August 2026}

\pagestyle{fancy}
\fancyhf{}
\fancyhead[L]{\footnotesize\itshape Coincidence Classes in the ECDSA State Space}
\fancyhead[R]{\footnotesize\itshape A. Polishchuk}
\fancyfoot[C]{\footnotesize\thepage}
\renewcommand{\headrulewidth}{0.4pt}
\fancypagestyle{plain}{\fancyhf{}\fancyfoot[C]{\footnotesize\thepage}%
  \renewcommand{\headrulewidth}{0pt}}

\begin{document}
\maketitle

\begin{abstract}
\noindent
We analyse a finite combinatorial object: the state space of an idealised ECDSA signer, together
with a family of subsets defined by exact arithmetic coincidences between the \emph{hidden} pair
$(\szk,\srxk) = (zk^{-1},\,rxk^{-1})$ and quantities computable from public data. Three such
subsets --- the classes $\clA$, $\clB$ and $\clE$ of a previously circulated
classification --- are treated in closed form.

The organising result is a conservation law. Let $R$ be \emph{any} rule mapping the public data of
a signature to a set $C$ of candidate values for the hidden component. Then
$\Pro[\szk \in C] = \Ex|C| / (n-2)$ exactly, so the success probability per tested candidate is the
constant $1/(n-2)$ for every rule, identical to sampling uniformly from $\Zs\setminus\{s\}$.
Classification cannot change the quality of candidates; it can change only their number. The
classes are subsets, their frequencies are cardinalities, and nothing further is available.

Instantiating the theorem yields exact frequencies, proved and then confirmed by exhaustive
enumeration of $1.66\times10^{9}$ states. Class $\clA$ emits $\Ex|C| = \Pro[\Qu] \to 1/4$
candidates per signature; class $\clE$ emits $5/6$ and is shown to be \emph{exactly} the affine
locus $\szk = 3\srxk + 2$, whence $\Pro[\clE] = 2\lfloor (n-3)/3\rfloor/(n-1)^{2} \sim 2/(3n)$;
class $\clB$ emits $\Theta(n^{-1}\log n)$, and its multiplier satisfies an \emph{integer}
quadratic, so no modular square root arises. Three structural facts follow: the qualification rate
is exactly $1/4$ on every curve, being the area of a lattice region; the ``Goldbach'' offset
distinguishing class $\clE$ is one of $\Theta(n)$ interchangeable choices; and the recovery
formulas of the classification are a single M\"obius bijection in disguise.

On secp256k1 the frequencies are $2^{-258.00}$, $2^{-256.58}$ and $\approx 2^{-506.9}$ per
signature. We report them because the question is natural, not because they are decisive: by the
main theorem the entire framework improves on uniform random key guessing by a factor of exactly
$1 + 1/(n-2) \approx 1 + 10^{-77}$, and that residual advantage is bought by the observation
$\szk \neq s$ rather than by any of the algebra. This is a negative result. Its value is that it
closes a natural line of enquiry with proofs rather than with estimates.
\end{abstract}

\vspace{4pt}
\noindent\textbf{Keywords:} ECDSA \textperiodcentered{} secp256k1 \textperiodcentered{}
lattice-point counting \textperiodcentered{} negative result \textperiodcentered{}
hidden number problem \textperiodcentered{} nonce distribution

\vspace{2pt}
\noindent\textbf{MSC 2020:} 11T71, 94A60, 11H06, 05A15

\vspace{6pt}
\hrule
\vspace{4pt}
\noindent\small\textbf{Statement of non-disclosure.} This paper reports no vulnerability. It
describes no attack, no partial attack and no weakening of any deployed parameter set. Its main
theorem proves that a family of proposed techniques cannot outperform exhaustive key search, which
is itself $\approx 2^{128}$ times more expensive than the generic attack the curve was designed to
resist. Readers arriving from the earlier ``vulnerability'' framing of this project's own
preprints are directed to Section~\ref{sec:errata}, in which that framing is withdrawn.
\vspace{4pt}
\hrule

\section{Introduction}

\subsection{Motivation}

ECDSA signing computes $s = (z + rx)k^{-1} \bmod n$, and the scalar $s$ splits additively into a
message part and a key part,
\[
  \szk = zk^{-1} \bmod n, \qquad \srxk = rxk^{-1}\bmod n, \qquad s \equiv \szk + \srxk \pmod n .
\]
Both summands are secret. An analyst with access to $(z,r,s)$ can nevertheless form auxiliary
integers --- most naturally $\szr = zr \bmod n$, the difference $d = \szr - s$, and the remainder
$a = s \bmod d$ --- and ask whether the hidden summands happen to stand in an exact arithmetic
relation to them. When they do, and when the relation pins $\szk$ to a small set of values, the
private key follows immediately, because $\szk$ determines the nonce and the nonce determines the
key.

Such coincidences do occur. On reduced-size test curves they occur often enough to be catalogued,
and a previously circulated framework for this project catalogued five of them, labelled $\clA$
through $\clE$, deriving closed-form key-recovery procedures for $\clA$, $\clB$ and $\clE$.
Empirically the classes thin out as the curve order $n$ grows, and the natural questions are how
fast, how far, and whether the thinning is a barrier or merely a nuisance.

This paper answers all three, and then shows that answering them was unnecessary.

\subsection{Contributions}

\begin{enumerate}[leftmargin=*,itemsep=3pt]
\item \textbf{A conservation law (Theorem~\ref{thm:zeroyield}).} For every
  \emph{abscissa-oblivious} rule --- one that treats $r$ as an opaque scalar and never invokes the
  group law --- the probability of a correct guess of $\szk$ equals the expected number of
  candidates emitted, divided by $n-2$. The yield per candidate is a constant, independent of the
  rule, so no such classification can beat uniform random guessing
  (Corollary~\ref{cor:nogo}); a rule that does beat it is obliged to solve ECDLP
  (Corollary~\ref{cor:ecdlp}). We claim no novelty for the underlying fact, which any
  cryptographer would regard as a restatement of what it means for a signature scheme to be
  correctly randomised. What the theorem contributes is a boundary: it converts an open-ended
  hunt for further coincidence classes into a closed question, and it makes precise which
  property a genuine improvement would have to violate.

\item \textbf{A structure theorem for class $\clE$ (Theorem~\ref{thm:Estruct}).} The three-clause
  ``Goldbach-type'' definition of class $\clE$ is equivalent to the single affine condition
  $\szk = 3\srxk + 2$ over $\Zz$. Its exact cardinality follows
  (Theorem~\ref{thm:Ecount}), as does the fact that the distinguished offset is arbitrary
  (Theorem~\ref{thm:offset}).

\item \textbf{Exact frequencies where the prior work had heuristics.} $\Pro[\Qu] = 1/4 +
  O(n^{-1}\log n)$ (Theorem~\ref{thm:qual}); $\Pro[\clA] = \Pro[\Qu]/(n-2)$, an identity
  (Theorem~\ref{thm:Afreq}); $\Pro[\clE] = 2\lfloor (n-3)/3\rfloor/(n-1)^{2}$, an identity
  (Theorem~\ref{thm:Ecount}); $\Pro[\clB] = \Theta(n^{-2}\log n)$
  (Proposition~\ref{prop:Bfreq}).

\item \textbf{Class $\clB$ lives over $\Zz$ (Theorem~\ref{thm:Bquad}).} Its multiplier satisfies
  an integer quadratic whose discriminant must be a perfect square --- a far stronger constraint
  than quadratic residuosity modulo $n$, and one that removes the modular square root from the
  published procedure. Wrap-around instances are shown to require $m \le 2$ together with an exact
  auxiliary relation (Lemma~\ref{lem:wrap}).

\item \textbf{Verification.} Exhaustive enumeration of the full state space for seven primes
  ($1.66\times10^{9}$ states), Monte-Carlo runs at $n \in \{9\,967,\ 99\,667\}$, and comparison
  against three independently generated test-curve datasets spanning five orders of magnitude in
  $n$. Every identity is confirmed to the precision of the enumeration; every asymptotic estimate
  is confirmed to within Poisson fluctuation.

\item \textbf{Errata (Section~\ref{sec:errata}).} Three corrections to previously circulated
  material, one of them to an earlier draft of this paper. The most consequential is that the
  published class-$\clE$ rate $1/(n-1)$ overstates the truth by $50\%$ and was already excluded at
  $32\sigma$ by the data accompanying it.
\end{enumerate}

\subsection{What is not a contribution}

Honesty about weight is part of the exercise. The following are \emph{not} contributions and are
not presented as such. The recovery formulas: Proposition~\ref{prop:master} shows they are one
M\"obius bijection. The taxonomy $\clA$--$\clE$ itself: an arbitrary partition of the analyst's
bookkeeping rather than of the object. The observation that rare coincidences occur at rate
$\Theta(1/n)$: folklore. And, as noted above, the substance of
Theorem~\ref{thm:zeroyield}, which is a formalisation rather than a discovery. The weight of this
work is that of a technical note recording a closed line of enquiry, and it is circulated as
such.

\subsection{What this paper does not claim}

It is not a vulnerability disclosure. It contains formulas that map a \emph{correct guess} of
$\szk$ to the private key, and those formulas are correct; but by Proposition~\ref{prop:master}
they constitute a single M\"obius bijection, that is, a change of coordinates on a search space of
size $n-1$, and by Theorem~\ref{thm:zeroyield} no public-data rule can make the guess better than
uniform. Reading them as an attack is a category error.

Nor, and this must be stated carefully, is the obstruction information-theoretic. ECDSA has no
information-theoretic security whatever: $r$ is the abscissa of $R = kG$, so it determines $R$ up
to the two points $\pm R$ (and, on secp256k1, an $x$-coordinate $r+n$ occurring with probability
$\approx 2^{-128}$), hence $k$ up to a handful of values, hence $x$. The conditional entropy of
the private key given one signature is a couple of bits, not $\log_2 n$. What makes ECDSA hard is
that extracting $k$ from $r$ is an instance of the elliptic-curve discrete logarithm problem.
Section~\ref{sec:zeroyield} therefore proves a statement about a restricted class of algorithms ---
those that treat $r$ as an opaque scalar and never invoke the group law --- and
Corollary~\ref{cor:ecdlp} makes the boundary explicit: any rule that beats the bound is obliged to
solve ECDLP.

\subsection{Related work}

Attacks that do recover ECDSA keys exploit a defect in the signer rather than a coincidence in the
output. Repeated nonces yield the key from two signatures by elementary algebra and have been
found repeatedly in the wild~\cite{bos2014,breitner2019}. Partially known or biased nonces reduce
to the Hidden Number Problem of Boneh and Venkatesan~\cite{boneh1996} and are solved by lattice
reduction~\cite{hgs2001,ns2002,ns2003}; a handful of leaked bits over a few dozen to a few million
signatures suffices, as demonstrated against real deployments by Breitner and
Heninger~\cite{breitner2019}, the Minerva group~\cite{minerva2020}, and
LadderLeak~\cite{ladderleak2020}, the last operating on less than one bit of leakage per
signature. Deterministic nonce derivation~\cite{rfc6979} removes the randomness failure mode at
the cost of introducing a fault-attack surface. The common feature is that each exploits
$\Theta(\ell)$ bits of information that is accessible by cheap algebra, whereas the classes studied
here supply none that an abscissa-oblivious rule can reach (Theorem~\ref{thm:zeroyield}). They
are therefore not a weaker point on the same spectrum but a different category. The information
that does remain in a signature --- $r$ determines $k$ up to sign --- is reachable only by solving
ECDLP, which is the province of~\cite{pollard1978} and not of any classification.

Generic attacks on the underlying discrete logarithm cost $\Theta(\sqrt{n}) \approx 2^{128}$
operations on secp256k1~\cite{pollard1978}; the curve parameters are those of SEC~2~\cite{sec2},
and the signature scheme is standardised in FIPS~186-5~\cite{fips1865}. The counting arguments
below use Dirichlet's divisor theorem in the form given by Hardy and Wright~\cite{hardywright},
and Landauer's principle~\cite{landauer1961} is quoted once, for scale only.

\subsection{Organisation}

Section~\ref{sec:prelim} fixes notation and proves the master identity.
Section~\ref{sec:model} states the probability model and its structure lemma.
Section~\ref{sec:zeroyield} proves the main theorem; a reader interested only in the conclusion
may stop there. Sections~\ref{sec:qual}--\ref{sec:E} instantiate it for the qualification event
and for classes $\clA$, $\clB$, $\clE$. Section~\ref{sec:numerics} reports the verification,
Section~\ref{sec:secp} evaluates on secp256k1,
Section~\ref{sec:errata} records the errata, and Section~\ref{sec:disc} discusses falsifiability
and limitations.

\section{Preliminaries}\label{sec:prelim}

\subsection{Signatures and hidden components}

Let $E/\mathbb{F}_p$ be an elliptic curve with base point $G$ of prime order $n$, and let $Q = xG$
be the public key of $x \in \{1,\dots,n-1\}$. For a message hash $z$ and nonce $k$, both in
$\{1,\dots,n-1\}$,
\begin{equation}\label{eq:sign}
  r = \bigl[kG\bigr]_x \bmod n, \qquad s = (z + rx)k^{-1}\bmod n .
\end{equation}
We write $\Zs = \{1,\dots,n-1\}$ and identify every element of $\mathbb{Z}_n$ with its
representative in $\{0,\dots,n-1\}$. This identification is essential, because the conditions below
mix modular arithmetic with ordinary integer comparison, divisibility and order. Relations meant
over $\Zz$ rather than over $\mathbb{Z}_n$ are marked as such.

\begin{definition}[Hidden components]\label{def:hidden}
$\zeta := \szk = zk^{-1}\bmod n$ and $\xi := \srxk = rxk^{-1}\bmod n$, so that
$s \equiv \zeta + \xi \pmod n$. Neither is observable from $(z,r,s)$.
\end{definition}

Since $\zeta,\xi \in \{1,\dots,n-1\}$, the \emph{integer} sum $\sigma := \zeta + \xi$ lies in
$\{2,\dots,2n-2\}$ and satisfies exactly one of
\begin{equation}\label{eq:wrap}
  \sigma = s \quad\text{(no wrap-around)}, \qquad\text{or}\qquad \sigma = s + n \quad\text{(wrap-around)}.
\end{equation}
The distinction is invisible publicly and matters for classes $\clB$ and $\clE$.

\begin{definition}[Public auxiliaries]\label{def:aux}
$\szr := zr \bmod n$; when $\szr > s$, $d := \szr - s$ and $a := s \bmod d$, the last an integer
remainder. All are computable from the signature and the signed message.
\end{definition}

\subsection{The master identity}

\begin{proposition}[Master identity; recovery is a relabelling]\label{prop:master}
\begin{enumerate}[label=(\roman*),leftmargin=2.2em,itemsep=2pt]
\item If $\hat{s} = \zeta$, then $x = z(s-\hat{s})(r\hat{s})^{-1}\bmod n$.
\item The map $\varphi : \Zs\setminus\{s\} \to \Zs$, $\varphi(\hat s) = z(s-\hat s)(r\hat s)^{-1}
  \bmod n$, is injective. Hence knowing $\zeta$ is equivalent to knowing $x$, and the formula in
  (i) transports no information from public data to the secret.
\end{enumerate}
\end{proposition}

\begin{proof}
(i) From Definition~\ref{def:hidden}, $k = z\zeta^{-1}$ and $\xi = (s - \zeta)\bmod n$.
Substituting into $\xi = rxk^{-1}$ gives
$x = \xi k r^{-1} = (s-\zeta)\cdot z\zeta^{-1}\cdot r^{-1} = z(s-\zeta)(r\zeta)^{-1}$.
(ii) Expand: $\varphi(\hat s) = zsr^{-1}\hat s^{-1} - zr^{-1}$. Inversion is a bijection of $\Zs$,
and $t \mapsto (zsr^{-1})t - zr^{-1}$ is an affine bijection since $zsr^{-1}\neq 0$ in the field
$\mathbb{Z}_n$. A composition of bijections is injective.
\end{proof}

Proposition~\ref{prop:master}(ii) governs the subject. Searching for $x$ through $\hat s$ is a
change of coordinates on a space of size $n-1$. Every class below is therefore no more than
\emph{a rule for proposing candidate values of $\zeta$}, and only two things about such a rule
matter: how often the proposal is correct, and how many proposals it costs.

\section{The probability model}\label{sec:model}

\begin{assumption}[Ideal signer]\label{as:U}
$z$, $x$ and $k$ are drawn independently and uniformly from $\Zs$.
\end{assumption}

\begin{assumption}[Abscissa heuristic]\label{as:R}
$r = [kG]_x \bmod n$ is modelled as uniform on $\Zs$ and independent of $(z,x,k)$.
\end{assumption}

Assumption~\ref{as:U} is the definition of a correctly implemented signer; when it fails, ECDSA
is broken by the mechanisms of Section~1.4, which have nothing to do with this paper.
Assumption~\ref{as:R} is a genuine heuristic, since $r$ is a deterministic function of $k$. It is
adopted only to obtain closed forms. Every prediction derived from it is checked in
Section~\ref{sec:numerics} against datasets generated on actual elliptic curves, where no such
assumption was made; agreement is within Poisson fluctuation across five orders of magnitude in
$n$ and eight in sample size. Assumption~\ref{as:R} is therefore not load-bearing for the
numerical conclusions.

\begin{lemma}[Structure of the state]\label{lem:structure}
Put $u := kr \bmod n$. Under Assumptions~\ref{as:U} and \ref{as:R}:
\begin{enumerate}[label=(\roman*),leftmargin=2.2em,itemsep=2pt]
\item $(\zeta,\xi,u)$ is uniform on $(\Zs)^{3}$;
\item the map $(\zeta,\xi,u)\mapsto(\zeta,s,\szr)$ with $s = \zeta+\xi$ and $\szr = \zeta u$
  (both mod $n$) is a bijection onto
  \[
    \mathcal{S} \;=\; \bigl\{(\zeta,s,\szr)\ :\ \zeta\in\Zs,\ s\in\mathbb{Z}_n\setminus\{\zeta\},\
    \szr\in\Zs\bigr\},\qquad |\mathcal{S}| = (n-1)^{3},
  \]
  so $(\zeta,s,\szr)$ is uniform on $\mathcal{S}$;
\item consequently $\szr$ is uniform on $\Zs$ and independent of $(\zeta,s)$, and for each fixed
  $s\neq 0$ the conditional law of $\zeta$ given $(s,\szr)$ is uniform on $\Zs\setminus\{s\}$, a
  set of exactly $n-2$ elements.
\end{enumerate}
\end{lemma}

\begin{proof}
(i) Fix $k,r$. Then $\zeta = zk^{-1}$ is a bijective image of the uniform $z$, hence uniform on
$\Zs$; $\xi = rxk^{-1}$ is a bijective image of the uniform $x$, hence uniform; they are
independent because $z \perp x$. Neither depends on $u = kr$, a function of $(k,r)$ alone; by
Assumption~\ref{as:R}, $u$ is uniform on $\Zs$ and independent of $(z,x)$, hence of $(\zeta,\xi)$.
(ii) The map is invertible by $\xi = s-\zeta$, $u = \szr\zeta^{-1}$, and its image is exactly
$\mathcal{S}$, since $\xi \neq 0 \iff s \neq \zeta$ and $u \neq 0 \iff \szr \neq 0$. A bijection
carries the uniform law to the uniform law, and $|\mathcal{S}| = (n-1)^{3}$ because for each
$\zeta$ there are $n-1$ admissible values of $s$.
(iii) Immediate from the product structure of $\mathcal{S}$.
\end{proof}

Lemma~\ref{lem:structure}(iii) is the workhorse. Within the model, after the analyst has computed
everything computable by modular arithmetic on $(z,r,s)$ --- $\szr$, $d$, $a$, and any function of
them --- the hidden component $\zeta$ remains uniform over $n-2$ possibilities.

The reader should be clear about the force of this. It is a statement about the model, not about
ECDSA. Assumption~\ref{as:R} is what makes it true, and Assumption~\ref{as:R} is false: $r$
determines $R = kG$ up to sign, hence $k$ up to a handful of candidates, hence $\zeta$ and $x$.
The residual entropy of $\zeta$ given a real signature is one or two bits. What
Assumption~\ref{as:R} does is sever precisely the link --- from the abscissa back to the scalar ---
whose difficulty is the entire security of the scheme. That severance is legitimate here for one
narrow reason, made precise in Definition~\ref{def:oblivious}: the classes under study never touch
the curve. It would be illegitimate for any conclusion about ECDSA in general, and no such
conclusion is drawn.

\section{The zero-yield theorem}\label{sec:zeroyield}

\begin{definition}[Public-data rule]\label{def:rule}
A \emph{rule} is a map $R$ from the public data of a signature --- $z$, $r$, $s$, and anything
computable from them, such as $\szr$, $d$, $a$ --- to a finite candidate set
$C = R(z,r,s) \subseteq \Zs$ for the hidden component $\zeta$. Randomised rules are permitted,
provided their randomness is independent of the signer's. The rule \emph{succeeds} if
$\zeta \in C$.
\end{definition}

\begin{definition}[Abscissa-oblivious rule]\label{def:oblivious}
A rule is \emph{abscissa-oblivious} if its output distribution is unchanged when the signer's
$r = [kG]_x \bmod n$ is replaced by an independent uniform draw from $\Zs$. Equivalently: the rule
treats $r$ as an opaque element of $\mathbb{Z}_n$ and never invokes the group law of $E$ --- it
never forms a curve point, never adds or doubles, never solves for a discrete logarithm.
\end{definition}

Classes $\clA$, $\clB$ and $\clE$ are abscissa-oblivious: every quantity they compute
($\szr = zr$, $d$, $a$, $m_1$, $m_2$, $\alpha$, $\beta$) is obtained by ring operations on
$z$, $r$, $s$ in $\mathbb{Z}_n$ together with integer division. The single scalar multiplication
each performs is a \emph{verification} of a finished candidate, not a step in producing one. The
theorem below applies to exactly this class of rules and to no wider one.

Every class of the classification is a rule. Class $\clA$ emits $C = \{a\}$ on qualified
transactions and $C = \varnothing$ otherwise. Class $\clE$ emits the roots of \eqref{eq:albe}
for each admissible branch of $\sigma$, discarding those outside $[1,n-1]$. Class $\clB$ emits
$\{am\}$ for each integer root $m$ of \eqref{eq:Bquad}. Uniform guessing is the rule that emits
$t$ i.i.d.\ uniform draws.

\begin{theorem}[Zero yield]\label{thm:zeroyield}
Under Assumptions~\ref{as:U} and~\ref{as:R} --- equivalently, for every abscissa-oblivious rule
$R$, for which Assumption~\ref{as:R} is by Definition~\ref{def:oblivious} no restriction ---
\begin{equation}\label{eq:zeroyield}
  \Pro\bigl[\zeta \in C\bigr]
  \;=\; \frac{\Ex\bigl|\,C \cap (\Zs\setminus\{s\})\,\bigr|}{n-2}
  \;\le\; \frac{\Ex|C|}{n-2}.
\end{equation}
In particular the success probability \emph{per candidate emitted} is the constant $1/(n-2)$,
independent of $R$.
\end{theorem}

\begin{proof}
By Lemma~\ref{lem:structure}(ii) the triple $(\zeta,s,\szr)$ is uniform on the product set
$\mathcal{S}$. Hence, conditionally on the public data --- a function of $(s,\szr,r)$, all
independent of $\zeta$ by Lemma~\ref{lem:structure}(iii) --- the hidden component $\zeta$ is
uniform on $\Zs\setminus\{s\}$, a set of exactly $n-2$ elements. The set $C$ is measurable with
respect to that same public data together with auxiliary randomness independent of the signer,
hence conditionally fixed. For a uniform variable on a set of size $n-2$ and a fixed subset $C$,
$\Pro[\zeta \in C \mid \text{public}] = |C \cap (\Zs\setminus\{s\})|/(n-2)$. Take expectations.
\end{proof}

\begin{corollary}[No-go for abscissa-oblivious rules]\label{cor:nogo}
No abscissa-oblivious classification of ECDSA signatures can achieve a per-candidate success rate
different from $1/(n-2)$. Consequently:
\begin{enumerate}[label=(\roman*),leftmargin=2.2em,itemsep=2pt]
\item no rule is \emph{better} than uniform guessing on $\Zs\setminus\{s\}$; equality in
  \eqref{eq:zeroyield} holds for that rule and for every rule emitting distinct candidates;
\item a rule can be \emph{worse} only by emitting duplicates or the excluded value $s$, both
  implementation defects rather than properties of the class;
\item against a guesser drawing uniformly from all of $\Zs$, whose yield is $1/(n-1)$, the best
  possible rule wins by the factor $(n-1)/(n-2) = 1 + 1/(n-2)$, equal to
  $1 + 8.64\times10^{-78}$ on secp256k1. That advantage is purchased entirely by the observation
  $\zeta \neq s$, which requires no algebra.
\end{enumerate}
\end{corollary}

\begin{corollary}[The boundary is ECDLP, not information]\label{cor:ecdlp}
A rule whose per-candidate yield exceeds $1/(n-2)$ must fail to be abscissa-oblivious: its output
must depend on the true $r = [kG]_x$ in a way that an independent uniform $r$ would not reproduce.
Any such dependence is a use of the group law of $E$. Since $r$ determines $R = \pm kG$ and the
step from $R$ to $k$ is by definition a discrete logarithm, such a rule solves an instance of
ECDLP.
\end{corollary}

\begin{proof}
Immediate from Definition~\ref{def:oblivious} and Theorem~\ref{thm:zeroyield}. If the rule's
output law is invariant under replacing $r$ by an independent uniform draw, the theorem applies
and the yield is $1/(n-2)$. Contrapositively, a yield above $1/(n-2)$ forces non-invariance, that
is, genuine dependence on the abscissa map $k \mapsto [kG]_x$.
\end{proof}

\begin{remark}[What is \emph{not} being claimed]\label{rem:notIT}
ECDSA offers no information-theoretic security, and nothing above suggests otherwise. A single
signature determines the private key up to a handful of candidates: $r$ fixes the abscissa of
$R = kG$, the curve equation gives $R$ up to sign, and $k$ follows, whence
$x = (sk-z)r^{-1}$. On secp256k1 the additional possibility $[R]_x = r+n$ arises with probability
$(p-n)/n \approx 2^{-128}$, so the candidate count is essentially two. The residual entropy of $x$
given $(z,r,s)$ is therefore one or two bits, not $256$ --- which is precisely why public-key
recovery from a signature is a routine operation. The hardness of ECDSA is computational and
resides entirely in ECDLP.

Assumption~\ref{as:R} deletes that structure. Doing so is what makes
Theorem~\ref{thm:zeroyield} provable, and it is equally what confines the theorem to
abscissa-oblivious rules. The theorem is worth stating because classes $\clA$, $\clB$, $\clE$ ---
and every classification of the same shape --- lie inside that confinement by construction. It
would be a serious error to read it as a claim that ECDSA signatures leak nothing.
\end{remark}

\begin{remark}\label{rem:framing}
The correct negative statement is an equality, not an inequality. ``Worse than guessing'' would
suggest an inefficiency that a cleverer rule might remove; ``exactly equal to guessing, for every
rule'' is a conservation law, and there is nothing left to optimise. Sections~\ref{sec:qual}
through \ref{sec:E} may therefore be read as a computation of the single scalar $\Ex|C|$ for three
particular rules, undertaken because the classes were already in circulation and their frequencies
had been misreported.
\end{remark}

\section{The qualification event}\label{sec:qual}

\begin{definition}[Qualification]\label{def:qual}
A transaction \emph{qualifies}, written $\Qu$, if $\szr > s$ and the residue $a = s \bmod d$
satisfies $a > 0$, $a \neq s$ and $a \neq \szr$.
\end{definition}

\begin{lemma}[Geometric form]\label{lem:qualgeom}
$\Qu$ holds if and only if $s < \szr \le 2s$ and $d \nmid s$.
\end{lemma}

\begin{proof}
Assume $\szr > s$, so $d \ge 1$. If $d > s$ then $s\bmod d = s$, violating $a \neq s$; conversely
if $d \le s$ then $a < d \le s$, so $a \neq s$. Thus $a \neq s \iff d \le s \iff \szr \le 2s$.
Next, $a = 0 \iff d \mid s$. Finally $a < d < \szr$ always, so $a \neq \szr$ is automatic.
\end{proof}

\begin{theorem}[Qualification rate]\label{thm:qual}
$\Pro[\Qu] = \tfrac14 - \Theta\!\left(n^{-1}\log n\right) \to \tfrac14$, independently of the
curve, the field and the key.
\end{theorem}

\begin{proof}
By Lemma~\ref{lem:structure}(iii) the pair $(s,\szr)$ is, up to an $O(1/n)$ perturbation of the
marginal of $s$, uniform on $\{1,\dots,n-1\}^{2}$. Write $M = n-1$, which is even for every odd
prime $n$, and count lattice points with $s < \szr \le \min(2s,M)$:
\[
  \sum_{s=1}^{M}\bigl(\min(2s,M) - s\bigr)
  = \underbrace{\sum_{s=1}^{M/2} s}_{=\,M^{2}/8 + M/4}
  \;+\; \underbrace{\sum_{s=M/2+1}^{M}(M-s)}_{=\,M^{2}/8 - M/4}
  \;=\; \frac{M^{2}}{4},
\]
an exact identity. The excluded set $d \mid s$ contributes
$\sum_{s \le M}\tau(s) = M\log M + (2\gamma-1)M + O(\sqrt{M})$ points by Dirichlet's divisor
theorem~\cite[Thm.~320]{hardywright}. Dividing by $M^{2}$ gives
$\Pro[\Qu] = \tfrac14 - (\log M + 2\gamma - 1)/M + O(M^{-1/2})$, and the perturbation of the
marginal of $s$ adds a further $O(1/M)$.
\end{proof}

\begin{remark}
The often-quoted ``$25\%$ anomaly rate'' of this project is thus the area of the region
$s < \szr \le 2s$ in the unit square. It is scale-invariant because that area is $1/4$ for every
$n$. It measures the geometry of comparing two independent residues and carries no information
about key recoverability; its constancy across curves is a reason to expect no exploitability, not
a reason to hope for it.
\end{remark}

\section{Class A: linear divisibility with unit multiplier}\label{sec:A}

\begin{definition}[Class $\clA$]\label{def:A}
A qualified transaction is in class $\clA$ if the ratio $m_1 = \zeta/a$, computed in $\mathbb{Q}$,
equals $1$; equivalently $\zeta = a$ as integers. (The strict taxonomy additionally demands
$m_2 = (s+\szr)/\xi \notin \Zz$, separating $\clA$ from $\clB$; by
Proposition~\ref{prop:Bfreq} that refinement removes an $O(n^{-2}\log n)$ fraction and is
invisible at every scale considered here.)
\end{definition}

\begin{theorem}[Recovery]\label{thm:Arec}
On a class-$\clA$ transaction, $k = az^{-1}\bmod n$ and
$x = (s-a)(rk)^{-1} = z(s-a)(ra)^{-1} \bmod n$. One candidate is produced, correct with certainty,
verified by one scalar multiplication $xG \stackrel{?}{=} Q$.
\end{theorem}

\begin{proof}
Substitute $\hat s = a = \zeta$ into Proposition~\ref{prop:master}(i). All inverses exist because
$n$ is prime and $z,r,a \not\equiv 0$.
\end{proof}

\begin{theorem}[Exact frequency]\label{thm:Afreq}
For every prime $n \ge 5$,
\begin{equation}\label{eq:Afreq}
  \Pro[\clA] \;=\; \frac{\Pro[\Qu]}{n-2}
  \;=\; \frac{1}{4(n-2)}\Bigl(1 - \Theta\bigl(n^{-1}\log n\bigr)\Bigr) \;\sim\; \frac{1}{4n},
\end{equation}
an identity in the model of Lemma~\ref{lem:structure}, not an asymptotic estimate. Equivalently,
the class-$\clA$ rule emits $\Ex|C| = \Pro[\Qu]$ candidates per signature.
\end{theorem}

\begin{proof}
Condition on $(s,\szr)$. If it does not qualify, $C = \varnothing$. If it qualifies then $d$ and
$a$ are determined, and by Lemma~\ref{lem:qualgeom} $1 \le a < d \le s$, so
$a \in \Zs\setminus\{s\}$ and $|C| = 1$. Apply Theorem~\ref{thm:zeroyield}.
\end{proof}

\begin{erratum}[A wrong turn worth recording]\label{err:divisor}
It is tempting to compute $\Pro[\clA]$ by counting, for each $(s,\zeta)$, the divisors
$d \mid (s-\zeta)$ with $d > \zeta$, and approximating the resulting sum by an integral. That
route yields $\Pro[\clA] \approx 0.3466/n$, which is $39\%$ too large and is rejected by the
tiny-curve data at $16.8\sigma$. The continuum approximation silently discards the fractional
parts $\{(n-\zeta)/d\}$, whose systematic mean of $\tfrac12$ accumulates to a term of order $d$
under the $\zeta$-summation. Carrying the floors exactly gives
$\sum_{\zeta=0}^{d-1}\lfloor (M-\zeta)/d\rfloor = M-d+1$ and restores $M^{2}/4$, in agreement with
Theorem~\ref{thm:Afreq}. Lattice counts in this subject are dominated by their floor terms;
integral approximations to them should be distrusted by default.
\end{erratum}

\section{Class B: dual integral divisibility}\label{sec:B}

\begin{definition}[Class $\clB$]\label{def:B}
A qualified transaction is in class $\clB$ if $m_1 = \zeta/a$ and $m_2 = (s+\szr)/\xi$ are both
integers and $m_1 = m_2 =: m$. Both quotients are taken in $\mathbb{Q}$: they assert exact
divisibility over $\Zz$, not congruences.
\end{definition}

\begin{theorem}[The multiplier satisfies an integer quadratic]\label{thm:Bquad}
On a class-$\clB$ transaction, with $\sigma \in \{s,\,s+n\}$ the true integer sum \eqref{eq:wrap},
\begin{equation}\label{eq:Bquad}
  a\,m^{2} \;-\; \sigma\,m \;+\; (s+\szr) \;=\; 0 \qquad\text{in } \Zz .
\end{equation}
Hence $\Delta = \sigma^{2} - 4a(s+\szr)$ is a perfect square in $\Zz$ and
$m = (\sigma \pm \sqrt{\Delta})/(2a) \in \Zz$, giving $\zeta = am$ and
$x = z(s-\zeta)(r\zeta)^{-1}\bmod n$. At most four candidates arise, two roots for each branch of
$\sigma$.
\end{theorem}

\begin{proof}
By definition $\zeta = am$ and $s + \szr = m\xi$ over $\Zz$. By \eqref{eq:wrap},
$\xi = \sigma - \zeta = \sigma - am$, so $s+\szr = m(\sigma - am) = \sigma m - am^{2}$, which is
\eqref{eq:Bquad}. All coefficients are integers, so the quadratic formula over $\Zz$ forces
$\Delta$ to be a perfect square.
\end{proof}

\begin{lemma}[Wrap-around in class $\clB$]\label{lem:wrap}
If a class-$\clB$ transaction has $\sigma = s+n$, then $m \le 2$; moreover $m = 1$ implies
$\szr + a = n$, and $m = 2$ implies that $d$ is even and $4a + d = 2n$.
\end{lemma}

\begin{proof}
Qualification gives $d \le s$, so $s + \szr = 2s + d \le 3s$. Also
$\xi = s + n - \zeta \ge s+1$ because $\zeta \le n-1$. From $m\xi = s+\szr$ we get
$m \le 3s/(s+1) < 3$. For $m=1$: $\xi = 2s+d$ and $\xi = s+n-a$ give $n = s+d+a = \szr + a$. For
$m=2$: $2\xi = 2s+d$ forces $d$ even with $\xi = s + d/2$, while $\xi = s+n-2a$; equating gives
$2n = 4a + d$.
\end{proof}

Each is one exact equation between quantities of size $\Theta(n)$, hence a codimension-one
condition on top of an already $O(n^{-2})$-rare event. No wrap-around instance occurred in the
exhaustive enumeration of Section~\ref{sec:numerics}. Implementations should nonetheless test both
branches; the cost is one extra scalar multiplication.

\begin{proposition}[Frequency of class $\clB$]\label{prop:Bfreq}
Under a divisor-equidistribution heuristic,
\begin{equation}\label{eq:Bfreq}
  \Pro[\clB] = \Theta\!\left(\frac{\log n}{n^{2}}\right),
  \qquad\text{empirically}\qquad
  \Pro[\clB] \approx \frac{c_B \log n}{n^{2}},\quad c_B \approx 0.19 .
\end{equation}
\end{proposition}

\begin{proof}[Argument]
Fix $(s,\szr)$ qualified. Each divisor $m$ of $s+\szr$ determines a unique candidate pair
$\zeta = am$, $\xi = (s+\szr)/m$; the pair survives only if additionally
$\zeta + \xi \in \{s,\,s+n\}$, a coincidence of probability $\Theta(1/n)$ under equidistribution.
The mean of $\tau$ on $[1,2n]$ is $\log n + O(1)$, so the conditional rate is
$\Theta(n^{-1}\log n)\cdot n^{-1}$; multiplying by $\Pro[\Qu] = \tfrac14$ gives
$\approx \tfrac14 n^{-2}\log n$. Exhaustive enumeration for $101 \le n \le 1009$ yields
$\Pro[\clB]\,n^{2}/\log n \in [0.189,\,0.196]$, so $c_B \approx 0.19$, within $25\%$ of the
heuristic constant $\tfrac14$.
\end{proof}

\begin{remark}[Weakest link, stated plainly]
The measured ratio drifts downward by about $3\%$ across the accessible range, so the true law may
carry a sub-logarithmic correction that the data cannot resolve; the accessible range spans one
decade of $n$ while the extrapolation target spans $74$. This is the least secure claim in the
paper. It is also inconsequential: an error of two orders of magnitude in $c_B$ moves
$\Pro[\clB]$ on secp256k1 from $10^{-153}$ to somewhere between $10^{-155}$ and $10^{-151}$, and
Theorem~\ref{thm:zeroyield} is unaffected either way. Classes $\clA$ and $\clE$ carry no such
caveat: their formulas are exact for every prime $n$.
\end{remark}

\section{Class E: the ``Goldbach-type'' decomposition}\label{sec:E}

Class $\clE$ is the most instructive of the three: its published description is ornate and its
actual content is one line.

\begin{definition}[Classes $\clE1$ and $\clE2$]\label{def:E}
With $\sigma = \zeta+\xi$ the integer sum \eqref{eq:wrap}, a transaction is in class $\clE$ if
$\sigma$ is even and
\begin{equation}\label{eq:Econd}
  \bigl|\zeta - \xi\bigr| \;=\; \frac{\sigma}{2} + 1 .
\end{equation}
Instances with $\sigma = s$ are labelled $\clE1$ and those with $\sigma = s+n$ are labelled
$\clE2$. The published formulation writes $\lfloor\sigma/2\rfloor + 1$; since $\sigma$ is required
to be even, the floor is vacuous.
\end{definition}

\begin{theorem}[Class $\clE$ is an affine line]\label{thm:Estruct}
For integers $\zeta,\xi \ge 1$, condition \eqref{eq:Econd} together with ``$\sigma$ even'' holds
if and only if
\begin{equation}\label{eq:Eline}
  \zeta = 3\xi + 2 \qquad\text{or}\qquad \xi = 3\zeta + 2 \qquad\text{in } \Zz .
\end{equation}
Moreover every class-$\clE$ instance satisfies $\sigma \equiv 2 \pmod 4$, and the two branches are
disjoint.
\end{theorem}

\begin{proof}
($\Rightarrow$) Equality \eqref{eq:Econd} rules out $\zeta = \xi$, its right-hand side being
positive. Take $\zeta > \xi$ without loss of generality; then
$\zeta - \xi = (\zeta+\xi)/2 + 1$, i.e.\ $2\zeta - 2\xi = \zeta + \xi + 2$, i.e.\
$\zeta = 3\xi+2$.
($\Leftarrow$) If $\zeta = 3\xi+2$ then $\sigma = 4\xi+2$ is even and
$\zeta - \xi = 2\xi + 2 = \sigma/2 + 1$, so \eqref{eq:Econd} holds; and $\sigma = 4\xi+2$ gives
$\sigma \equiv 2 \pmod 4$. If both branches held, $\zeta = 3(3\zeta+2)+2 = 9\zeta+8$ forces
$\zeta = -1$, impossible for a positive integer.
\end{proof}

\begin{remark}[The Goldbach framing is empty]
Theorem~\ref{thm:Estruct} dissolves the name. The published description --- ``a symmetric pair
around $\sigma/2$ with fixed difference'' --- is doubly vacuous: \emph{every} pair summing to
$\sigma$ is symmetric about $\sigma/2$ by construction, and the only content of \eqref{eq:Econd}
is that the hidden pair lies on one particular line of slope $3$. There is no additive number
theory here and no connection to the Goldbach problem beyond the superficial fact that a number is
written as a sum of two others.
\end{remark}

\begin{theorem}[Exact cardinality and frequency]\label{thm:Ecount}
For every prime $n \ge 5$ the number of class-$\clE$ hidden pairs in $\{1,\dots,n-1\}^{2}$ is
\begin{equation}\label{eq:Ecount}
  \#\clE \;=\; 2\left\lfloor \frac{n-3}{3}\right\rfloor,
  \qquad\text{so}\qquad
  \Pro[\clE] \;=\; \frac{2\lfloor (n-3)/3\rfloor}{(n-1)^{2}} \;=\; \frac{2}{3n} + O\!\left(\frac{1}{n^{2}}\right).
\end{equation}
\end{theorem}

\begin{proof}
By Theorem~\ref{thm:Estruct} the class is a disjoint union of two branches. On $\zeta = 3\xi+2$
the constraints are $\xi \ge 1$ and $3\xi+2 \le n-1$, i.e.\ $1 \le \xi \le (n-3)/3$, giving
$\lfloor (n-3)/3\rfloor$ solutions; the other branch gives the same count by symmetry. Class
$\clE$ depends only on $(\zeta,\xi)$, uniform on $(\Zs)^{2}$ by
Lemma~\ref{lem:structure}(i), a set of $(n-1)^{2}$ elements. No class-$\clE$ pair is removed by
$s \neq 0$, since $\sigma = 4\xi+2 = n$ would force $n$ even.
\end{proof}

\begin{theorem}[Recovery]\label{thm:Erec}
Given the effective sum $\sigma \in \{s,\,s+n\}$ with $\sigma \equiv 2 \pmod 4$, put
\begin{equation}\label{eq:albe}
  \alpha = \frac{3\sigma+2}{4}, \qquad \beta = \sigma - \alpha = \frac{\sigma-2}{4}.
\end{equation}
Then $\{\zeta,\xi\} = \{\alpha,\beta\}$, and $x$ is one of
$x_c = z(s-c)(rc)^{-1}\bmod n$ for $c \in \{\alpha,\beta\}$, disambiguated by one scalar
multiplication.
\end{theorem}

\begin{proof}
Solve $\zeta+\xi = \sigma$, $\zeta-\xi = \sigma/2+1$: adding gives $2\zeta = 3\sigma/2 + 1$, i.e.\
$\zeta = (3\sigma+2)/4$, an integer precisely because $\sigma \equiv 2 \pmod 4$
(Theorem~\ref{thm:Estruct}). Apply Proposition~\ref{prop:master}(i) to each candidate.
\end{proof}

\begin{theorem}[Offset invariance]\label{thm:offset}
For any fixed $c \in \Zz$ define $\clE_c$ by ``$\sigma$ even and
$|\zeta-\xi| = \sigma/2 + c$''. Then
$\clE_c = \{\zeta = 3\xi+2c\} \sqcup \{\xi = 3\zeta+2c\}$, and for $|c| = o(n)$,
\[
  \Pro[\clE_c] = \frac{2}{3n}\Bigl(1 + O\bigl(|c|/n\bigr)\Bigr),
\]
independently of $c$. The Goldbach value $c=1$ is thus one member of a family of $\Theta(n)$
equally frequent classes.
\end{theorem}

\begin{proof}
Repeat Theorem~\ref{thm:Estruct} with right-hand side $\sigma/2 + c$: for $\zeta > \xi$,
$2\zeta - 2\xi = \zeta+\xi+2c$ gives $\zeta = 3\xi+2c$; conversely $\sigma = 4\xi+2c$ is even and
$\zeta - \xi = 2\xi+2c = \sigma/2 + c$. Counting on one branch,
$\#\{\xi : 1 \le \xi \le n-1,\ 1 \le 3\xi+2c \le n-1\} = n/3 + O(|c|)$; the other is symmetric.
\end{proof}

A control experiment in the project repository reports a flat frequency of $\approx 0.67/n$ for
offsets $c \in \{\pm1,3,5,7,9,101\}$. Theorem~\ref{thm:offset} accounts for it exactly: the
constant $0.67 \approx 2/3$ is $\Pro[\clE_c]\cdot n$, and its independence of $c$ is a two-line
consequence of the affine structure. An experiment run to test a hypothesis turns out to be a
corollary of that hypothesis's own algebra.

\subsection{Class versus detector}

The rule of Definition~\ref{def:rule} associated with class $\clE$ must try both branches of
$\sigma$, since it cannot see which is real. Counting admissible candidates:
$\sigma = s$ contributes two whenever $s \equiv 2 \pmod 4$; $\sigma = s+n$ contributes two when
additionally $s \le n/3$ and one otherwise, because $\alpha$ then exceeds $n-1$ while $\beta$ does
not. Hence
\begin{equation}\label{eq:EC}
  \Ex|C_{\clE}| = \tfrac14\cdot 2 + \tfrac14\Bigl(\tfrac13\cdot 2 + \tfrac23\cdot 1\Bigr) = \tfrac56,
  \qquad \Pro[\text{rule } \clE \text{ succeeds}] = \frac{5}{6(n-2)} ,
\end{equation}
which exceeds $\Pro[\clE] = 2/(3n)$ of Theorem~\ref{thm:Ecount}. There is no contradiction: the
\emph{class} is defined using the secret $\sigma$, whereas the \emph{rule} also scores accidental
hits from the wrong branch, and the gap $5/6$ against $2/3$ is exactly their weight. This
distinction --- between a class defined with oracle access to the secret and a rule that must
operate without it --- is where informal treatments of the subject usually go wrong.

\section{Numerical validation}\label{sec:numerics}

\subsection{Exhaustive enumeration}

For small $n$ the entire state space $\mathcal{S}$ of Lemma~\ref{lem:structure}, of size
$(n-1)^{3}$, can be enumerated without sampling, giving exact probabilities. Table~\ref{tab:exh}
reports all states for seven primes, $1.66\times10^{9}$ in total.

\begin{table}[t]
\centering\small
\begin{tabular}{rrrrrrrr}
\toprule
$n$ & states & $\Pro[\Qu]$ & $\Pro[\clA]$ & $\Pro[\Qu]/(n-2)$ & $\#\clE$ & $2\lfloor\frac{n-3}{3}\rfloor$ & $\Pro[\clB]n^{2}/\ln n$\\
\midrule
101 & 1.00e6 & 0.2096820 & 2.118000e{-}3 & 2.118000e{-}3 & 64 & 64 & 0.192\\
151 & 3.38e6 & 0.2205200 & 1.480000e{-}3 & 1.480000e{-}3 & 98 & 98 & 0.195\\
251 & 1.56e7 & 0.2303389 & 9.250560e{-}4 & 9.250560e{-}4 & 164 & 164 & 0.196\\
401 & 6.40e7 & 0.2364823 & 5.926875e{-}4 & 5.926875e{-}4 & 264 & 264 & 0.193\\
503 & 1.26e8 & 0.2388253 & 4.766967e{-}4 & 4.766967e{-}4 & 332 & 332 & 0.194\\
751 & 4.22e8 & 0.2419628 & 3.230483e{-}4 & 3.230483e{-}4 & 498 & 498 & 0.191\\
1009 & 1.03e9 & 0.2437219 & 2.420277e{-}4 & 2.420277e{-}4 & 670 & 670 & 0.189\\
\bottomrule
\end{tabular}
\caption{Exhaustive enumeration of the full state space. Columns 4--5 confirm
Theorem~\ref{thm:Afreq} as an identity to every displayed digit; columns 6--7 confirm
\eqref{eq:Ecount} exactly; the last column calibrates $c_B \approx 0.19$ in \eqref{eq:Bfreq}.
$\Pro[\Qu]$ approaches $1/4$ from below at the rate $\Theta(\log n/n)$ of
Theorem~\ref{thm:qual}.}
\label{tab:exh}
\end{table}

\subsection{Independently generated test-curve data}

The datasets accompanying the earlier framework were generated on genuine elliptic curves, so
Assumption~\ref{as:R} played no role in producing them. Table~\ref{tab:curves} compares those
counts with the formulas proved above. All residuals lie within Poisson fluctuation; the largest,
$+2.05\sigma$ on $12$ observed events, has a one-sided $p$-value of $0.05$. Measured qualification
rates were $0.2494$ and $0.2503$ on the two million-transaction datasets, against $1/4$ predicted.

\begin{table}[t]
\centering\small
\begin{tabular}{lrrrrrrrrrr}
\toprule
curve & $n$ & $N_{\mathrm{tx}}$ & A pred & A obs & $z_{\clA}$ & B pred & B obs & E pred & E obs & $z_{\clE}$\\
\midrule
tiny  & 9\,967 & 9.93e7 & 2\,491.8 & 2\,465 & $-0.54$ & 1.75 & 2 & 6\,642.2 & 6\,746 & $+1.27$\\
small & 99\,667 & 1.00e6 & 2.51 & 4 & $+0.94$ & 2e{-}4 & 0 & 6.69 & 12 & $+2.05$\\
large & 1\,241\,630\,743 & 1.00e6 & 2e{-}4 & 0 & $-0.01$ & 3e{-}12 & 0 & 5e{-}4 & 0 & $-0.02$\\
\bottomrule
\end{tabular}
\caption{Predicted versus observed counts on the project's test curves.}
\label{tab:curves}
\end{table}

\subsection{Monte-Carlo cross-check and the yield constant}

Simulating the model of Lemma~\ref{lem:structure} with $2.4\times10^{7}$ draws per curve gave, at
$n = 9\,967$: qualification rate $0.24909$; $\#\clA = 600$ against $599.8$ predicted;
$\#\clE = 1\,614$ against $1\,605.3$. At $n = 99\,667$: qualification $0.24978$, $\#\clA = 71$
(predicted $60.1$), $\#\clE = 177$ (predicted $160.5$). In all $4.8\times10^{7}$ draws the number
of transactions satisfying Definition~\ref{def:E} was \emph{identical} to the number satisfying
the affine criterion \eqref{eq:Eline} --- $1\,614 = 1\,614$ and $177 = 177$ --- confirming
Theorem~\ref{thm:Estruct} without exception.

Table~\ref{tab:yield} tests Theorem~\ref{thm:zeroyield} directly by measuring, for each rule, both
the expected candidate count and the empirical yield per candidate.

\begin{table}[t]
\centering\small
\begin{tabular}{lrrrr}
\toprule
rule & $\Ex|C|$ predicted & $\Ex|C|$ measured & yield/candidate measured & $1/(n-2)$\\
\midrule
class $\clA$ & $\Pro[\Qu] \to 1/4$ & 0.2489 & $1.031\times10^{-4}$ & $1.004\times10^{-4}$\\
class $\clE$ & $5/6$ & 0.8325 & $9.809\times10^{-5}$ & $1.004\times10^{-4}$\\
class $\clB$ & $\approx 0.19\ln n/n$ & $0.184\ln n/n$ & --- & $1.004\times10^{-4}$\\
\bottomrule
\end{tabular}
\caption{Direct test of Theorem~\ref{thm:zeroyield} at $n = 9\,967$ over $3\times10^{6}$
simulated signatures. The yield per candidate is the same constant for every rule, as the theorem
requires; class $\clB$ emitted too few candidates at this sample size for a yield estimate.}
\label{tab:yield}
\end{table}

\section{Evaluation on secp256k1}\label{sec:secp}

The group order is the prime
\[
  n = \texttt{FFFFFFFF FFFFFFFF FFFFFFFF FFFFFFFE}\ \texttt{BAAEDCE6 AF48A03B BFD25E8C D0364141}_{16}
\]
\[
  \phantom{n}\approx 1.157920892\times10^{77},
\]
with $\log_2 n = 256.000000$, $\ln n = 177.4457$, $n \equiv 1 \pmod 3$ and $n \equiv 1 \pmod 4$.
Because $n \equiv 1 \pmod 3$ the floor in \eqref{eq:Ecount} evaluates exactly to $(n-4)/3$, so no
approximation enters the class-$\clA$ or class-$\clE$ figures.

\begin{table}[t]
\centering\small
\begin{tabular}{llrll}
\toprule
event & closed form & value & power of $2$ & status\\
\midrule
$\Qu$ & $\tfrac14 + O(n^{-1}\log n)$ & $0.250000\ldots$ & $2^{-2.000}$ & theorem\\
$\clA$ & $\Pro[\Qu]/(n-2)$ & 2.159042e{-}78 & $2^{-258.00}$ & exact identity\\
$\clB$ & $\approx 0.19\ln n/n^{2}$ & $\approx$2.5e{-}153 & $\approx 2^{-506.9}$ & calibrated heuristic\\
$\clE$ & $2\lfloor\frac{n-3}{3}\rfloor/(n-1)^{2}$ & 5.757446e{-}78 & $2^{-256.58}$ & exact identity\\
$\clA\cup\clB\cup\clE$ & union bound & 7.916488e{-}78 & $2^{-256.24}$ & ---\\
\emph{one blind guess} & $1/(n-1)$ & 8.636169e{-}78 & $2^{-256.00}$ & exact\\
\bottomrule
\end{tabular}
\caption{Per-signature probabilities on secp256k1. These are per-signature figures; the
like-for-like comparison at equal numbers of tested candidates is Corollary~\ref{cor:nogo}.}
\label{tab:secp}
\end{table}

Taking $N_{\mathrm{tx}} = 10^{10}$ as a generous upper bound on all secp256k1 signatures ever
produced --- Bitcoin has processed on the order of $10^{9}$ transactions carrying a few
$\times 10^{9}$ input signatures, Ethereum a few $\times10^{9}$ transactions --- the expected
number of instances of all three classes combined is $7.9\times10^{-68}$; at $N_{\mathrm{tx}} =
10^{12}$ it is $7.9\times10^{-66}$, and at an absurd $10^{20}$ it is $7.9\times10^{-58}$. Since
these expectations are far below $1$, Markov's inequality makes them upper bounds on the
probability that even one such signature exists.

The sample size for even odds of a single occurrence is $N_{1/2} = \ln 2/\Pro$: $1.20\times10^{77}$
for $\clE$, $3.21\times10^{77}$ for $\clA$, $8.76\times10^{76}$ for the union and
$2.76\times10^{152}$ for $\clB$. Every entry exceeds the group order.

\paragraph{Abundance without density.}
The classes are not empty on secp256k1. The transaction space
$\mathcal{T} = \{1,\dots,n-1\}^{4}$ has $|\mathcal{T}| \approx 1.80\times10^{308}$ elements, of
which $\#\clA \approx 3.88\times10^{230}$, $\#\clE \approx 1.04\times10^{231}$ and
$\#\clB \approx 4.5\times10^{155}$. These are enormous absolute counts and a fraction
$\approx 10^{-77}$ of the space. An analyst who reasons ``there are $10^{231}$ vulnerable
signatures out there'' has said something true and entirely without consequence, in the same way
that there are $10^{231}$ private keys beginning with any particular $256$-bit prefix. Absolute
counts in this subject should never be quoted without the density beside them.

\begin{remark}[Physical bound, for scale only]
Landauer's principle~\cite{landauer1961} charges at least $k_BT\ln 2 = 2.61\times10^{-23}\,$J per
irreversible bit operation at the cosmic microwave background temperature. Even at one bit erasure
per trial, enumerating $1/\Pro[\clA] = 4.63\times10^{77}$ trials costs $1.21\times10^{55}\,$J, the
mass--energy of $6.8\times10^{7}$ solar masses. The figure is quoted once, for scale, and no
argument in this paper rests on it: Corollary~\ref{cor:nogo} already settles the matter at the
level of algorithms, where it belongs.
\end{remark}

\section{Errata to previously circulated material}\label{sec:errata}

This project circulated earlier write-ups of classes $\clA$ and $\clB$ under a ``vulnerability''
framing, together with a framework document stating frequency formulas for all five classes. That
framing is withdrawn, and three specific numerical or structural claims are corrected here.

\begin{erratum}[The class-$\clE$ rate is $2/(3n)$, not $1/(n-1)$]\label{err:E}
The framework document states $\Pro[\clE] = 1/(n-1)$, reasoning that the condition ``fixes $\szk$
to at most two values'' and that ``exactly one of $s$, $s+n$ is even''. The second clause is
correct; the first over-counts. By Theorem~\ref{thm:Estruct} an even $\sigma$ admits a solution
only when $\sigma \equiv 2 \pmod 4$, which removes half the even sums, and the surviving locus has
length $\approx n/3$ per branch rather than $\approx n/2$. The correct rate is $2/(3n)$, that is,
$\tfrac23$ of the published value.

The distinction is not academic. With $n = 9\,967$ and $N_{\mathrm{tx}} = 99\,323\,922$ --- the
project's own tiny-curve dataset, published alongside the formula --- the model $1/(n-1)$ predicts
$9\,966.5$ events against $6\,746$ observed, a Poisson deviation of $-32.3\sigma$;
Theorem~\ref{thm:Ecount} predicts $6\,642.2$, a deviation of $+1.27\sigma$. The published formula
was excluded by the data accompanying it.
\end{erratum}

\begin{erratum}[No modular square root is required in class $\clB$]\label{err:B}
The published statement of the class-$\clB$ recovery reduces the quadratic modulo $n$, computes
$\sqrt{\Delta} \bmod n$ via the Euler criterion, and inverts $2a$ modulo $n$. That procedure is
not incorrect, but it is strictly weaker and needlessly expensive: it discards the fact that
\eqref{eq:Bquad} holds over $\Zz$, where ``$\Delta$ is a perfect square'' is far more restrictive
than ``$\Delta$ is a quadratic residue mod $n$'' (density $\Theta(\Delta^{-1/2})$ against $1/2$).
It also fixes $\sigma = s$, silently assuming no wrap-around; see Lemma~\ref{lem:wrap}. In
exhaustive enumeration for $n \le 1009$ --- $3\,823$ class-$\clB$ instances --- the integer
discriminant was a perfect square in every case.
\end{erratum}

\begin{erratum}[Correction to an earlier draft of this paper]\label{err:self}
An earlier draft compared an ``anomaly hunter'' processing $T$ signatures against a guesser making
$3T$ guesses, concluding that the hunter loses by a factor of $3.3$. The arithmetic is right but
the accounting is arbitrary: it charges the hunter for three candidates per signature, whereas the
classes emit $\Ex|C| \approx \tfrac14 + \tfrac56 + O(n^{-1}\log n) \approx 1.083$. Charging both
adversaries the same number of scalar multiplications gives an \emph{exact tie}, up to the factor
$1 + 1/(n-2)$ of Corollary~\ref{cor:nogo}. The corrected statement is stronger, for the reason
given in Remark~\ref{rem:framing}.
\end{erratum}

\section{Discussion}\label{sec:disc}

\subsection{Falsifiability}

The results are conditional on Assumptions~\ref{as:U} and~\ref{as:R}. The useful question is what
observation would overturn them. A class whose observed frequency exceeds its predicted $c/n$
baseline by a factor $\lambda > 1$ that persists as $n$ grows would demonstrate a genuine
non-uniformity in the law of $\zeta$ --- the only finding in this area worth pursuing, since it
would contradict Lemma~\ref{lem:structure} and hence Theorem~\ref{thm:zeroyield}.

Such a test is cheap on test curves. For a Poisson count with baseline rate $p_0$, distinguishing
$p_0$ from $\lambda p_0$ at significance $\alpha$ and power $1-\beta$ requires
\begin{equation}\label{eq:power}
  N \;\gtrsim\; \frac{\bigl(z_{\alpha/2} + z_{\beta}\bigr)^{2}}{p_0\bigl(\sqrt{\lambda}-1\bigr)^{2}} .
\end{equation}
At $\alpha = 0.05$ and $80\%$ power, detecting $\lambda = 2$ in class $\clE$ on the tiny curve
needs $6.8\times10^{5}$ transactions; $\lambda = 1.10$ needs $4.9\times10^{7}$; $\lambda = 1.01$
needs $4.7\times10^{9}$. The existing $9.9\times10^{7}$-transaction dataset already resolves a
$10\%$ enrichment.

Applying \eqref{eq:power} to the data in hand: the tiny-curve class-$\clE$ count ($6\,746$ against
$6\,642.2$ predicted) constrains $\lambda$ to $1.016 \pm 0.012$, excluding any enrichment above
$4\%$ at $95\%$ confidence; the class-$\clA$ count ($2\,465$ against $2\,491.8$) constrains
$\lambda$ to $0.989 \pm 0.020$. Both are consistent with $\lambda = 1$ to within a couple of
percent.

\begin{remark}[The one experiment worth running]
Searching for further coincidence classes is not it: the space of arithmetic coincidences is
infinite and, by Theorem~\ref{thm:zeroyield}, each new class will have frequency
$\Ex|C|/(n-2)$ with $\Ex|C| = \Theta(1)$ at best. The informative experiment is a direct test of
the \emph{distribution}: take $10^{8}$ to $10^{9}$ transactions on a curve with
$n \approx 10^{4}$--$10^{5}$, form the empirical law of $\zeta = zk^{-1}$ conditioned on every
public statistic available, and test uniformity by a chi-squared statistic over $\sqrt{N}$ bins or
by a discrete Fourier test in the manner of Bleichenbacher. A null result closes the question; a
positive one would announce itself as a bias in a histogram, not as an elegant identity.
\end{remark}

\subsection{Limitations}

\begin{enumerate}[leftmargin=*,itemsep=3pt]
\item \textbf{Assumption~\ref{as:R} is a heuristic.} Treating $r$ as an independent uniform unit
  is not provable with current techniques. Its role is confined to obtaining closed forms, and the
  predictions were checked against data generated without it (Table~\ref{tab:curves}).
\item \textbf{The class-$\clB$ constant is empirical.} $c_B \approx 0.19$ is calibrated over
  $101 \le n \le 1009$ and drifts by $\approx 3\%$ across that range. Only the order
  $\Theta(n^{-2}\log n)$ is argued from first principles, and that rests on divisor
  equidistribution.
\item \textbf{Ideal signer only.} Nothing here bears on implementations with faulty randomness,
  deterministic-nonce bugs, fault injection or side channels; those are the real attack surface
  and are governed by the works cited in Section~1.4.
\item \textbf{Union bounds.} The three classes share the same hidden pair and are not independent,
  so the union figures in Table~\ref{tab:secp} are upper bounds --- the conservative direction.
\item \textbf{Classes $\clC$ and $\clD$ are out of scope.} Class $\clC$ admits no closed-form
  recovery. Class $\clD$ is not a coincidence class but the generic case, occurring for
  $\approx 25\%$ of transactions by Theorem~\ref{thm:qual}; its candidate window has width
  $a = \Theta(n)$ on average, so $\Ex|C_{\clD}| = \Theta(n)$ and
  Theorem~\ref{thm:zeroyield} gives success probability $\Theta(1)$ at cost $\Theta(n)$ --- that
  is, exhaustive search, roughly $2^{128}$ times worse than Pollard rho. The level index
  $N = \lfloor \zeta/a\rfloor$ is moreover a function of the secret $\zeta$, so the window cannot
  be located without already knowing what it is meant to reveal.
\end{enumerate}

\section{Conclusion}

Under an ideal signer the hidden pair is uniform on $(\Zs)^{2}$ and, jointly with the public
$\szr$, uniform on a product set. Conditioning on all public data therefore leaves $\zeta$ uniform
over $n-2$ values, and every subsequent result is a consequence of that one sentence: the yield
per tested candidate is the constant $1/(n-2)$ for every rule computable from public data, so a
classification can change only how many candidates are emitted, never their quality. The best
conceivable rule of this type beats uniform random key guessing by the factor
$1 + 8.6\times10^{-78}$, and exhaustive key search is itself $\approx 2^{128}$ times more
expensive than the generic attack secp256k1 was designed to resist.

Within that frame the three classes are exercises. Class $\clA$ is the graph of a function and
occurs with probability $\Pro[\Qu]/(n-2)$; class $\clE$ is the affine line
$\zeta = 3\xi+2$ and occurs with probability $2\lfloor (n-3)/3\rfloor/(n-1)^{2}$; class $\clB$ is
a divisor coincidence of order $n^{-2}\log n$ whose multiplier satisfies an integer quadratic. The
qualification rate is the area of a lattice region and equals $1/4$ on every curve. The
distinguished offset of class $\clE$ is one of $\Theta(n)$ equivalent choices.

The productive question is not whether more classes exist --- infinitely many do, all of the same
quality --- but whether the distribution of $\szk$ is uniform. That question is answerable with
existing data and existing tools, and the answer so far, to within $\pm 2\%$, is yes.

\subsection*{Data and code availability}

All figures are reproduced by the enumeration in Appendix~\ref{app:code}, which is exhaustive for
$n \le 1009$ and therefore yields exact rationals rather than estimates. Datasets, the generator,
and the control experiment are in the project repository at \url{[repository URL]}.

\begin{thebibliography}{99}\small

\bibitem{ladderleak2020}
D.~F. Aranha, F.~R. Novaes, A.~Takahashi, M.~Tibouchi, and Y.~Yarom.
\newblock LadderLeak: Breaking ECDSA with less than one bit of nonce leakage.
\newblock In \emph{ACM CCS 2020}, pages 225--242. ACM, 2020.
\newblock DOI 10.1145/3372297.3417268; ePrint 2020/615.

\bibitem{boneh1996}
D.~Boneh and R.~Venkatesan.
\newblock Hardness of computing the most significant bits of secret keys in Diffie--Hellman and
related schemes.
\newblock In \emph{CRYPTO '96}, LNCS 1109, pages 129--142. Springer, 1996.

\bibitem{bos2014}
J.~W. Bos, J.~A. Halderman, N.~Heninger, J.~Moore, M.~Naehrig, and E.~Wustrow.
\newblock Elliptic curve cryptography in practice.
\newblock In \emph{Financial Cryptography 2014}, LNCS 8437, pages 157--175. Springer, 2014.

\bibitem{breitner2019}
J.~Breitner and N.~Heninger.
\newblock Biased nonce sense: Lattice attacks against weak ECDSA signatures in cryptocurrencies.
\newblock In \emph{Financial Cryptography 2019}, LNCS 11598, pages 3--20. Springer, 2019.
\newblock ePrint 2019/023.

\bibitem{sec2}
Certicom Research.
\newblock SEC 2: Recommended elliptic curve domain parameters, version 2.0, 2010.
\newblock \url{https://www.secg.org/sec2-v2.pdf}.

\bibitem{hardywright}
G.~H. Hardy and E.~M. Wright.
\newblock \emph{An Introduction to the Theory of Numbers}.
\newblock Oxford University Press, 6th edition, 2008.

\bibitem{hgs2001}
N.~A. Howgrave-Graham and N.~P. Smart.
\newblock Lattice attacks on digital signature schemes.
\newblock \emph{Designs, Codes and Cryptography}, 23(3):283--290, 2001.

\bibitem{minerva2020}
J.~Jancar, V.~Sedlacek, P.~Svenda, and M.~Sys.
\newblock Minerva: The curse of ECDSA nonces.
\newblock \emph{IACR Transactions on Cryptographic Hardware and Embedded Systems},
2020(4):281--308, 2020.

\bibitem{landauer1961}
R.~Landauer.
\newblock Irreversibility and heat generation in the computing process.
\newblock \emph{IBM Journal of Research and Development}, 5(3):183--191, 1961.

\bibitem{ns2002}
P.~Q. Nguyen and I.~E. Shparlinski.
\newblock The insecurity of the digital signature algorithm with partially known nonces.
\newblock \emph{Journal of Cryptology}, 15(3):151--176, 2002.

\bibitem{ns2003}
P.~Q. Nguyen and I.~E. Shparlinski.
\newblock The insecurity of the elliptic curve digital signature algorithm with partially known
nonces.
\newblock \emph{Designs, Codes and Cryptography}, 30(2):201--217, 2003.

\bibitem{fips1865}
National Institute of Standards and Technology.
\newblock \emph{Digital Signature Standard (DSS)}. FIPS PUB 186-5, 2023.

\bibitem{pollard1978}
J.~M. Pollard.
\newblock Monte Carlo methods for index computation $(\mathrm{mod}\ p)$.
\newblock \emph{Mathematics of Computation}, 32(143):918--924, 1978.

\bibitem{rfc6979}
T.~Pornin.
\newblock Deterministic usage of the Digital Signature Algorithm (DSA) and Elliptic Curve Digital
Signature Algorithm (ECDSA).
\newblock RFC 6979, IETF, 2013.

\end{thebibliography}

\appendix

\section{Reproduction code}\label{app:code}

The following is self-contained. For $n \le 1009$ it enumerates all $(n-1)^{3}$ states, so the
probabilities it returns are exact. Every assertion below is a restatement of a theorem in the
body, and all of them pass.

\begin{lstlisting}
import numpy as np, math

def enumerate_states(n):
    """Exhaustive sweep over the state space S of the structure lemma: (zeta, xi, u)."""
    v = np.arange(1, n, dtype=np.int64)
    Z, X = np.meshgrid(v, v, indexing='ij')
    Z, X = Z.ravel(), X.ravel()
    S = (Z + X) % n
    keep = S != 0
    Z, X, S = Z[keep], X[keep], S[keep]

    sig = Z + X                                       # true integer sum
    nE  = int(np.count_nonzero((sig % 2 == 0) &
                               (np.abs(Z - X) == sig // 2 + 1)))
    nE_line = int(np.count_nonzero((Z == 3*X + 2) | (X == 3*Z + 2)))  # E structure thm
    assert nE == nE_line == 2 * ((n - 3) // 3)                        # exact E count

    q = a_ct = b_ct = 0
    for u in range(1, n):
        ZR = (Z * u) % n
        m  = ZR > S
        z_, x_, s_, zr_ = Z[m], X[m], S[m], ZR[m]
        d  = zr_ - s_
        a  = s_ % d
        g  = (a > 0) & (a != s_) & (a != zr_)                         # qualification
        z_, x_, s_, zr_, a = z_[g], x_[g], s_[g], zr_[g], a[g]
        q += a.size
        a_ct += int(np.count_nonzero(z_ == a))                        # class A
        t = s_ + zr_
        both = (z_ % a == 0) & (t % x_ == 0)                          # class B
        if both.any():
            b_ct += int(np.count_nonzero(
                z_[both] // a[both] == t[both] // x_[both]))

    tot = (n - 1) ** 3
    assert abs(a_ct/tot - (q/tot)/(n-2)) < 1e-15                      # exact A frequency
    return dict(n=n, PQ=q/tot, PA=a_ct/tot, PE=nE/(n-1)**2, PB=b_ct/tot)

for n in (101, 151, 251, 401, 503, 751, 1009):
    r = enumerate_states(n)
    print(f"n={r['n']:5d}  Pr[Q]={r['PQ']:.6f}  Pr[A]={r['PA']:.6e}  "
          f"Pr[E]={r['PE']:.6e}  Pr[B]*n^2/ln n={r['PB']*n*n/math.log(n):.3f}")

# --- secp256k1 -------------------------------------------------------------
from decimal import Decimal as D, getcontext
getcontext().prec = 60
n  = 0xFFFFFFFFFFFFFFFFFFFFFFFFFFFFFFFEBAAEDCE6AF48A03BBFD25E8CD0364141
PA = D(1) / 4 / (n - 2)
PE = D(2 * ((n - 3) // 3)) / D((n - 1) ** 2)
PB = D("0.19") * D(n).ln() / (D(n) * D(n))
print("Pr[A] = %.6E   Pr[E] = %.6E   Pr[B] = %.3E   1/n = %.6E"
      % (PA, PE, PB, D(1) / n))
\end{lstlisting}

\section{Notation}\label{app:notation}

\noindent
\begin{tabular}{@{}ll@{}}
$n$ & prime group order\\
$\zeta = \szk = zk^{-1}$, \ $\xi = \srxk = rxk^{-1}$ & hidden components (Definition~\ref{def:hidden})\\
$\sigma = \zeta + \xi \in \{s,\,s+n\}$ & integer sum \eqref{eq:wrap}\\
$\szr = zr$, \ $d = \szr - s$, \ $a = s \bmod d$ & public auxiliaries (Definition~\ref{def:aux})\\
$m_1 = \zeta/a$, \ $m_2 = (s+\szr)/\xi$ & divisibility ratios, taken in $\mathbb{Q}$\\
$\Qu$ & qualification event (Definition~\ref{def:qual})\\
$C$, \ $\Ex|C|$ & candidate set of a rule; its expected size\\
$\tau(\cdot)$, \ $\gamma$ & divisor function; Euler--Mascheroni constant\\
\end{tabular}

\medskip\noindent
All modular arithmetic is modulo $n$ unless a relation is explicitly stated to hold in $\Zz$.

\end{document}
